\setchoices{label=(\Alph*)}

\begin{question}[index=1]
已知函数 $f(x)=\dfrac{\displaystyle\int_0^x|\sin t|\dd t}{x^\alpha}$ 在 $(0,+\infty)$ 内有界，则 $\alpha$ 的取值范围为
\begin{choices}[columns=4]
\item $[0,2]$
\item $[1,2]$
\item $(1,3]$
\item $[1,3)$
\end{choices}
\end{question}
\QuestionBreak

\begin{question}[index=2]
设积分 $\displaystyle\int_L\frac{ax}{y^3}\dd x+\frac{y^2-3x^2}{y^4}\dd y$ 在上半平面 $(y>0)$ 上与路径无关，则 $\dfrac{ax}{y^3}\dd x+\dfrac{y^2-3x^2}{y^4}\dd y$ \allowbreak\mbox{的一个原函数} $u(x,y)=$
\begin{choices}[columns=2]
\item $y^2+\dfrac{x^3}{y^3}-\dfrac1y$
\item $x^2+\dfrac{x^3}{y^3}-\dfrac1y$
\item $\dfrac{x^2}{y^3}-\dfrac1y$
\item $\dfrac{x^2}{y}+\dfrac{x^3}{y^3}-\dfrac1y$
\end{choices}
\end{question}
\QuestionBreak

\begin{question}[index=3]
级数 $\displaystyle\sum_{n=1}^{\infty}\ln(1+u_n+u_n^2)$ 收敛，则下列选项一定正确的是
\begin{choices}[columns=2]
\item $\displaystyle\sum_{n=1}^{\infty}(-1)^nu_n$ 收敛
\item $\displaystyle\sum_{n=1}^{\infty}u_n^2$ 收敛
\item $\displaystyle\sum_{n=1}^{\infty}u_n$ 发散
\item $\displaystyle\sum_{n=1}^{\infty}\frac{u_n}{n^2}$ 绝对收敛
\end{choices}
\end{question}
\QuestionBreak

\begin{question}[index=4]
设 $\displaystyle M=\int_0^1\eu^x\dd x,\quad N=\int_0^{\pi/2}\eu^{\sin x}\dd x,\quad P=\int_0^{\pi/4}\eu^{\tan x}\dd x$，则
\begin{choices}[columns=4]
\item $P<M<N$
\item $P<N<M$
\item $M<P<N$
\item $N<P<M$
\end{choices}
\end{question}
\QuestionBreak

\begin{question}[index=5]
设 $\vect A=(\vect\alpha_1,\vect\alpha_2,\vect\alpha_3)$，其中 $\vect\alpha_1,\vect\alpha_2,\vect\alpha_3$ 为两两正交的三维非零列向量，且
\[
\vect B=(\vect\alpha_1-\vect\alpha_2,\vect\alpha_2+\lambda\vect\alpha_3,\vect\alpha_3-2\vect\alpha_1)\text{。}
\]
若 $|\vect A|=|\vect B|$，则实数 $\lambda=$
\begin{choices}[columns=4]
\item $-1$
\item $0$
\item $\dfrac12$
\item $2$
\end{choices}
\end{question}
\QuestionBreak

\begin{question}[index=6]
设 $\vect A,\vect B,\vect C$ 均为 $n$ 阶方阵，且 $\vect A$ 可逆，则下列矩阵与 $\mat{\vect A&\vect B\\\vect O&\vect C}$ 的秩可能不相等的为
\begin{choices}[columns=4]
\item $\mat{\vect A&\vect O\\\vect O&\vect C}$
\item $\mat{\vect E&\vect O\\\vect A&\vect C}$
\item $\mat{\vect A&\vect B\vect C\\\vect O&\vect C}$
\item $\mat{\vect A&\vect O\\\vect O&\vect B\vect C}$
\end{choices}
\end{question}
\QuestionBreak

\begin{question}[index=7]
设 $\vect P$ 为 $3$ 阶可逆矩阵，且 $\vect P^{\mathrm T}=\vect P^*$，$\vect P^*$ 为 $\vect P$ 的伴随矩阵，$\vect A$ 为 $3$ 阶实对称矩阵，\allowbreak\mbox{$\vect x$ 为三维列向量}，二次型 $\vect x^{\mathrm T}\vect A\vect x$ 经变换 $\vect x=\vect P\vect y$ 化为二次型 $\vect y^{\mathrm T}\vect B\vect y$，则 $\vect A,\vect B$ 的关系为
\begin{choices}[columns=2]
\item 等价、相似、且合同
\item 等价、相似、未必合同
\item 等价、合同、未必相似
\item 合同、相似、未必等价
\end{choices}
\end{question}
\QuestionBreak

\begin{question}[index=8]
甲和乙进行乒乓球比赛，假设每局比赛甲胜的概率为 $\dfrac23$，采取三局两胜制，若最终甲获胜，则比赛共进行了三局的概率为
\begin{choices}[columns=4]
\item $\dfrac25$
\item $\dfrac35$
\item $\dfrac12$
\item $\dfrac23$
\end{choices}
\end{question}
\QuestionBreak

\begin{question}[index=9]
设随机变量 $X$ 的概率密度函数为 $f(x)$，期望方差均存在，其中 $D(X)=\sigma^2$，随机变量 $Y$ 的\allowbreak\mbox{概率密度函数}为 $f(-y)$，且 $X$ 与 $Y$ 的相关系数为 $\dfrac12$，记 $Z=X+Y$，则根据\allowbreak\mbox{切比雪夫不等式}得 $P\{|Z|\geq 3\sigma\}\leq$
\begin{choices}[columns=4]
\item $\dfrac13$
\item $\dfrac19$
\item $\dfrac23$
\item $\dfrac29$
\end{choices}
\end{question}
\QuestionBreak

\begin{question}[index=10]
设 $\overline X$ 为来自总体 $X\dist N(\mu,\sigma^2)$（$\sigma^2$ 已知）的一个容量为 $n$ 的简单随机样本的样本均值，\allowbreak\mbox{若已知}在置信水平 $1-2\alpha$ 下 $\mu$ 的置信区间的长度为 $2$，则在显著性水平 $\alpha$ 下，对于\allowbreak\mbox{假设检验问题} $H_0:\mu\geq 1,\ H_1:\mu<1$，若检验结果拒绝 $H_0$，则应有
\begin{choices}[columns=4]
\item $\overline x\in(-\infty,1)$
\item $\overline x\in(-\infty,-1)$
\item $\overline x\in(-\infty,0)$
\item $\overline x\in(0,+\infty)$
\end{choices}
\end{question}
